\documentclass[12pt,a4paper]{article}

\usepackage[utf8]{inputenc}
\usepackage[T1]{fontenc}
\usepackage{geometry}
\usepackage{graphicx}
\usepackage{hyperref}
\usepackage{tocloft}
\usepackage{fancyhdr}
\usepackage{titlesec}
\usepackage{listings}
\usepackage{xcolor}
\usepackage{amsmath}
\usepackage{algorithm}
\usepackage{algpseudocode}
\usepackage{booktabs}
\usepackage{caption}
\usepackage{subcaption}
\usepackage{tcolorbox}
\usepackage{dirtree}
% \usepackage{algorithmic}

\geometry{
    left=1in,
    right=1in,
    top=1in,
    bottom=1in
}

\hypersetup{
    colorlinks=true,
    linkcolor=blue,
    filecolor=magenta,
    urlcolor=cyan,
    citecolor=green,
    pdftitle={VectorDB: High-Performance Vector Database System},
    pdfauthor={},
    pdfsubject={Data Structures and Algorithms},
    pdfkeywords={Vector Database, Embeddings, Approximate Nearest Neighbor, Graph Algorithms, Data Structures}
}

\pagestyle{fancy}
\fancyhf{}
\fancyhead[L]{\leftmark}
\fancyhead[R]{\thepage}
\renewcommand{\headrulewidth}{0.4pt}

\titleformat{\section}
{\normalfont\Large\bfseries}{\thesection}{1em}{}
\titleformat{\subsection}
{\normalfont\large\bfseries}{\thesubsection}{1em}{}
\titleformat{\subsubsection}
{\normalfont\normalsize\bfseries}{\thesubsubsection}{1em}{}

\lstdefinestyle{cpp}{
    language=C++,
    basicstyle=\ttfamily\footnotesize,
    keywordstyle=\color{blue}\bfseries,
    commentstyle=\color{gray}\itshape,
    stringstyle=\color{red},
    numbers=left,
    numberstyle=\tiny\color{gray},
    stepnumber=1,
    numbersep=10pt,
    backgroundcolor=\color{gray!10},
    showspaces=false,
    showstringspaces=false,
    showtabs=false,
    frame=single,
    tabsize=4,
    captionpos=b,
    breaklines=true,
    breakatwhitespace=false,
    escapeinside={\%*}{*)},
    morekeywords={Vector, VectorDB, SearchEngine, GraphANN, Heap, MinHeap, MaxHeap}
}

\lstset{style=cpp}

\begin{document}

\begin{titlepage}
    \centering
    \vspace*{1cm}
    
    {\Huge\bfseries Vector Database Implementation\\[0.3cm]}
    
    \vspace{1cm}
    
    \includegraphics[width=0.35\textwidth]{image.png}
    
    \vspace{1cm}
    
    {\Large\bfseries Session 2024 -- 2028}
    
    \vspace{1.5cm}
    
    {\Large\bfseries Submitted by:}\\[0.5cm]
    \begin{tabular}{cc}
        {\large Hussain Shahbaz} & {\large Ali Hassan Siddiqui} \\
        {\large 2024-CS-04} & {\large 2024-CS-05}
    \end{tabular}
    
    \vspace{1.5cm}
    
    {\Large\bfseries Supervised by:} {\large Mr. Nazeef Ul Haq}\\[0.3cm]
    {\Large\bfseries Course:}\\[0.2cm]
    {\large CSC-202 DSA}
    
    \vspace{1.5cm}
    
    {\Large\bfseries Department of Computer Science}\\[0.2cm]
    {\large University of Engineering and Technology}\\[0.2cm]
    {\large Lahore Pakistan}
    
\end{titlepage}

\newpage
\tableofcontents
\newpage

\listoffigures
\listoftables

\newpage

\begin{abstract}
This document presents VectorDB, a high-performance vector database system implemented entirely from scratch using custom data structures in C++. The system enables efficient storage, indexing, and retrieval of high-dimensional vector embeddings generated from textual data. VectorDB addresses the growing need for semantic search capabilities in modern applications by implementing state-of-the-art algorithms including hierarchical navigable small world graphs (HNSW-inspired), brute-force search, and persistent storage with compression. The project demonstrates advanced understanding of fundamental data structures including dynamic arrays, heaps, graphs, B-trees, hash maps, and Huffman encoding. The system integrates with Google's Gemini API for generating 128-dimensional embeddings and provides both exact and approximate nearest neighbor search capabilities. Performance analysis shows that the graph-based approximate search achieves significant speedup compared to brute-force methods while maintaining high recall rates. This implementation serves as both a practical tool for semantic search applications and an educational resource demonstrating real-world applications of classical data structures and algorithms.
\end{abstract}

\newpage

\section{Introduction}

\subsection{Background}
Vector databases have emerged as a critical infrastructure component in modern artificial intelligence and machine learning systems. Unlike traditional databases that store structured data in rows and columns, vector databases specialize in storing and retrieving high-dimensional vectors that represent semantic meanings of data objects such as text, images, or audio.

The proliferation of large language models and embedding models has created an urgent need for efficient vector storage and retrieval systems. Applications ranging from semantic search engines to recommendation systems rely heavily on the ability to find similar vectors quickly among millions or billions of candidates.

\subsection{Motivation}
Traditional database systems are inadequate for semantic similarity search because:
\begin{itemize}
    \item High-dimensional vectors cannot be efficiently indexed using conventional B-tree or hash-based indices
    \item Similarity metrics (cosine similarity, Euclidean distance) require specialized algorithms
    \item Exact nearest neighbor search in high dimensions suffers from the curse of dimensionality
    \item Modern applications demand sub-millisecond query latency even with millions of vectors
\end{itemize}

This project was motivated by the desire to understand and implement the fundamental algorithms and data structures that power commercial vector databases such as Pinecone, Weaviate, and Milvus.

\subsection{Project Scope}
VectorDB is a comprehensive vector database system that includes:
\begin{itemize}
    \item Custom implementation of all data structures (no STL containers used)
    \item Integration with Gemini API for generating text embeddings
    \item Multiple search strategies: brute-force and graph-based approximate nearest neighbor
    \item Persistent storage with B-tree indexing and Huffman compression
    \item Comprehensive testing suite for validation and performance analysis
\end{itemize}

\section{Problem Statement}

\subsection{Core Problem}
Given a collection of $n$ high-dimensional vectors $\mathcal{V} = \{v_1, v_2, \ldots, v_n\}$ where each $v_i \in \mathbb{R}^d$ (typically $d = 128, 384, 768$), and a query vector $q \in \mathbb{R}^d$, efficiently find the $k$ vectors in $\mathcal{V}$ that are most similar to $q$ according to a similarity metric $\text{sim}(v_i, q)$.

\subsection{Challenges}

\subsubsection{Computational Complexity}
The naive approach of computing similarity between the query and all stored vectors has time complexity $O(nd)$, which becomes prohibitive as $n$ grows to millions or billions.

\subsubsection{Memory Efficiency}
Storing $n$ vectors of dimension $d$ with 32-bit floats requires $4nd$ bytes. For $n = 10^6$ and $d = 128$, this equals approximately 512 MB, necessitating efficient memory management and compression.

\subsubsection{Dimensionality Curse}
Traditional spatial indexing structures (k-d trees, R-trees) degrade to linear search performance in high dimensions, requiring novel indexing approaches.

\subsubsection{Accuracy-Speed Tradeoff}
Approximate search algorithms must balance between search speed and result quality (recall), requiring careful parameter tuning.

\subsection{Requirements}

\subsubsection{Functional Requirements}
\begin{enumerate}
    \item \textbf{Vector Insertion}: Support insertion of vectors with unique identifiers
    \item \textbf{Exact Search}: Implement brute-force k-nearest neighbor search
    \item \textbf{Approximate Search}: Implement graph-based ANN for fast retrieval
    \item \textbf{Persistence}: Save and load vector database from disk
    \item \textbf{API Integration}: Connect to embedding generation service
    \item \textbf{Compression}: Reduce storage footprint using Huffman encoding
\end{enumerate}

\subsubsection{Non-Functional Requirements}
\begin{enumerate}
    \item \textbf{Performance}: Sub-second query time for databases with $10^3$ to $10^5$ vectors
    \item \textbf{Scalability}: Handle databases growing to millions of vectors
    \item \textbf{Accuracy}: Maintain $>$ 90\% recall@10 for approximate search
    \item \textbf{Code Quality}: Clean, modular, well-tested implementation
    \item \textbf{Educational Value}: Demonstrate practical application of DSA concepts
\end{enumerate}

\section{System Architecture}

\subsection{Overall Design}
VectorDB follows a modular architecture with clear separation of concerns:

\begin{figure}[h]
    \centering
    \begin{verbatim}
    ┌─────────────────────────────────────────┐
    │         User Interface Layer            │
    │    (CLI / API for queries & inserts)    │
    └────────────────┬────────────────────────┘
                     │
    ┌────────────────▼────────────────────────┐
    │         Search Engine Layer             │
    │  (Brute Force / Graph ANN Search)       │
    └────────────────┬────────────────────────┘
                     │
    ┌────────────────▼────────────────────────┐
    │      Vector Storage Layer (VectorDB)    │
    │    (In-memory vector management)        │
    └─────┬──────────────────────────┬────────┘
          │                          │
    ┌─────▼──────────┐     ┌────────▼─────────┐
    │  Embedding API │     │  Disk Storage    │
    │  (Gemini API)  │     │  (B-Tree + Huff) │
    └────────────────┘     └──────────────────┘
    \end{verbatim}
    \caption{VectorDB System Architecture}
\end{figure}

\subsection{Module Descriptions}

\subsubsection{Data Structures Layer}
All fundamental data structures implemented from scratch:
\begin{itemize}
    \item \textbf{Vector}: Dynamic array with amortized $O(1)$ append
    \item \textbf{Heap}: Priority queue for top-k maintenance
    \item \textbf{Graph}: Adjacency list for HNSW graph
    \item \textbf{HashMap}: Hash table for ID lookups
    \item \textbf{BTree}: Disk-based index for persistent storage
    \item \textbf{Huffman}: Compression tree for vector encoding
    \item \textbf{List}: Linked list for dynamic collections
    \item \textbf{Queue/Stack}: FIFO/LIFO structures for algorithms
    \item \textbf{Pair}: Key-value tuple structure
\end{itemize}

\subsubsection{Core Vector Layer}
\begin{itemize}
    \item \textbf{VectorItem}: Encapsulates vector ID and embedding
    \item \textbf{VectorDB}: Main database class managing vector collection
    \item \textbf{Similarity}: Computes cosine and Euclidean metrics
    \item \textbf{SearchEngine}: Implements brute-force k-NN search
    \item \textbf{GraphANN}: Implements approximate search using graphs
    \item \textbf{KDANN}: Implements KD-Tree based nearest neighbor search
    \item \textbf{BallANN}: Implements Ball Tree based nearest neighbor search
\end{itemize}

\subsubsection{Embedding Layer}
\begin{itemize}
    \item \textbf{Embedding}: Interfaces with Gemini API for text embedding
    \item Generates 128-dimensional vectors from natural language input
    \item Handles API authentication and response parsing
\end{itemize}

\subsubsection{Storage Layer}
\begin{itemize}
    \item \textbf{B-Tree Index}: Maps vector IDs to disk offsets
    \item \textbf{Huffman Compression}: Reduces storage requirements
    \item Persistent serialization and deserialization
\end{itemize}

\section{Data Structures Implementation}

\subsection{Dynamic Vector (Array)}

\subsubsection{Design}
A dynamically resizable array that grows by doubling capacity when full, providing amortized $O(1)$ insertion time.

\subsubsection{Operations}
\begin{itemize}
    \item \texttt{push\_back(T)}: $O(1)$ amortized
    \item \texttt{operator[](int)}: $O(1)$
    \item \texttt{pop\_back()}: $O(1)$
    \item \texttt{resize(int)}: $O(n)$
\end{itemize}

\subsubsection{Memory Management}
Uses manual memory allocation with proper copy constructor and destructor to prevent memory leaks.

\subsubsection{Usage in VectorDB}
\begin{itemize}
    \item Stores array of \texttt{VectorItem*} pointers
    \item Stores float embeddings within each \texttt{VectorItem}
    \item Stores search results (vector of pairs)
\end{itemize}

\subsection{Min/Max Heap}

\subsubsection{Design}
Template-based heap implementation supporting both min and max heap operations through virtual comparison function.

\subsubsection{Heapify Algorithm}
\begin{algorithm}
\caption{Heapify Up}
\begin{algorithmic}[1]
\State $index \gets \text{last position}$
\While{$\text{hasParent}(index)$ \textbf{and} $\text{compare}(\text{heap}[index], \text{heap}[\text{parent}(index)])$}
    \State $\text{swap}(\text{heap}[index], \text{heap}[\text{parent}(index)])$
    \State $index \gets \text{parent}(index)$
\EndWhile
\end{algorithmic}
\end{algorithm}

\subsubsection{Usage in VectorDB}
\begin{itemize}
    \item \textbf{MinHeap}: Maintains top-k largest similarities during search
    \item \textbf{MaxHeap}: Can be used for bottom-k queries
    \item Stores \texttt{Pair<float, int>} for (similarity, vectorID)
\end{itemize}

\subsubsection{Complexity Analysis}
\begin{itemize}
    \item Insert: $O(\log n)$
    \item Extract-min/max: $O(\log n)$
    \item Peek: $O(1)$
    \item Build heap: $O(n)$
\end{itemize}

\subsection{Graph (Adjacency List)}

\subsubsection{Design}
Directed graph using adjacency list representation where each node stores a list of neighbor IDs.

\subsubsection{Operations}
\begin{itemize}
    \item \texttt{addNode()}: $O(1)$
    \item \texttt{addEdge(u, v)}: $O(1)$
    \item \texttt{getNeighbors(u)}: $O(1)$ access, $O(degree(u))$ traversal
    \item \texttt{hasEdge(u, v)}: $O(degree(u))$
\end{itemize}

\subsubsection{Usage in VectorDB}
\begin{itemize}
    \item Represents HNSW-inspired navigation graph
    \item Each vector is a node
    \item Edges connect vectors to their k-nearest neighbors
    \item Enables fast greedy search through graph traversal
\end{itemize}

\subsection{HashMap}

\subsubsection{Design}
Template-based hash table with separate chaining for collision resolution.

\subsubsection{Hash Function}
Uses modulo-based hashing with prime table sizes to minimize collisions.

\subsubsection{Usage in VectorDB}
\begin{itemize}
    \item ID to VectorItem mapping for $O(1)$ retrieval
    \item Caching API responses
    \item Tracking visited nodes during graph search
\end{itemize}

\subsubsection{Load Factor Management}
Automatically resizes when load factor exceeds threshold (typically 0.75) to maintain performance.

\subsection{B-Tree}

\subsubsection{Design}
Self-balancing tree structure optimized for disk I/O with large branching factor.

\subsubsection{Properties}
\begin{itemize}
    \item Each node contains multiple keys (order $m$)
    \item All leaf nodes at same level
    \item Minimum occupancy: $\lceil m/2 \rceil - 1$ keys per node
    \item Maximum occupancy: $m - 1$ keys per node
\end{itemize}

\subsubsection{Usage in VectorDB}
\begin{itemize}
    \item Maps vector IDs to file offsets on disk
    \item Enables efficient range queries
    \item Optimizes disk reads by minimizing seeks
\end{itemize}

\subsubsection{Complexity}
\begin{itemize}
    \item Search: $O(\log_m n)$
    \item Insert: $O(\log_m n)$
    \item Delete: $O(\log_m n)$
\end{itemize}

\subsection{Huffman Tree}

\subsubsection{Design}
Binary tree for optimal prefix-free encoding based on symbol frequencies.

\subsubsection{Encoding Algorithm}
\begin{algorithm}
\caption{Huffman Encoding}
\begin{algorithmic}[1]
\State Count frequency of each symbol
\State Create leaf node for each symbol with frequency
\State Build min-heap of nodes
\While{heap size $> 1$}
    \State $left \gets \text{extractMin}()$
    \State $right \gets \text{extractMin}()$
    \State Create internal node with frequency = $left.freq + right.freq$
    \State $internal.left \gets left$
    \State $internal.right \gets right$
    \State Insert $internal$ into heap
\EndWhile
\State Root = remaining node in heap
\end{algorithmic}
\end{algorithm}

\subsubsection{Usage in VectorDB}
\begin{itemize}
    \item Compresses float vectors before writing to disk
    \item Reduces storage by 30-50\% depending on data distribution
    \item Decompresses vectors when loading from disk
\end{itemize}

\subsection{Auxiliary Structures}

\subsubsection{List (Linked List)}
\begin{itemize}
    \item Singly linked list with head and tail pointers
    \item Used in graph adjacency lists
    \item Dynamic insertion/deletion without reallocation
\end{itemize}

\subsubsection{Stack}
\begin{itemize}
    \item LIFO structure used in graph DFS traversal
    \item Used in Huffman tree construction
\end{itemize}

\subsubsection{Queue}
\begin{itemize}
    \item FIFO structure for BFS traversal
    \item Level-order processing in trees
\end{itemize}

\subsubsection{Pair}
\begin{itemize}
    \item Generic key-value tuple
    \item Used extensively: (similarity, ID), (ID, offset)
    \item Supports comparison operators for heap operations
\end{itemize}

\section{Core Algorithms}

\subsection{Similarity Metrics}

\subsubsection{Cosine Similarity}
Measures the cosine of the angle between two vectors:
\begin{equation}
\text{cosine}(u, v) = \frac{u \cdot v}{||u|| \cdot ||v||} = \frac{\sum_{i=1}^{d} u_i v_i}{\sqrt{\sum_{i=1}^{d} u_i^2} \cdot \sqrt{\sum_{i=1}^{d} v_i^2}}
\end{equation}

Range: $[-1, 1]$ where 1 indicates identical direction, -1 opposite, 0 orthogonal.

\textbf{Complexity}: $O(d)$ where $d$ is dimension.

\subsubsection{Euclidean Distance}
Measures straight-line distance in vector space:
\begin{equation}
\text{euclidean}(u, v) = \sqrt{\sum_{i=1}^{d} (u_i - v_i)^2}
\end{equation}

Range: $[0, \infty)$ where smaller values indicate greater similarity.

\textbf{Complexity}: $O(d)$

\subsection{Brute Force Search}

\subsubsection{Algorithm}
\begin{algorithm}
\caption{Brute Force k-NN Search}
\begin{algorithmic}[1]
\Require Query vector $q$, database $\mathcal{D}$, $k$
\Ensure Top-$k$ most similar vectors
\State Initialize MinHeap $H$ of size $k$
\For{each vector $v_i$ in $\mathcal{D}$}
    \State $sim \gets \text{similarity}(q, v_i)$
    \If{$|H| < k$}
        \State $H.\text{insert}((sim, i))$
    \ElsIf{$sim > H.\text{min}().similarity$}
        \State $H.\text{extractMin}()$
        \State $H.\text{insert}((sim, i))$
    \EndIf
\EndFor
\State Extract all elements from $H$ and reverse order
\State \Return sorted results
\end{algorithmic}
\end{algorithm}

\subsubsection{Complexity Analysis}
\begin{itemize}
    \item Time: $O(nd + n \log k)$ where $n$ = database size, $d$ = dimension
    \item Space: $O(k)$ for heap
    \item Guarantees exact results (100\% recall)
\end{itemize}

\subsubsection{Optimization}
\begin{itemize}
    \item Early termination: Stop if remaining vectors cannot improve results
    \item SIMD vectorization: Parallelize similarity computation
    \item Batch processing: Compute multiple queries together
\end{itemize}

\subsection{Graph-Based Approximate Search (HNSW-Inspired)}

\subsubsection{Graph Construction}
\begin{algorithm}
\caption{Build Navigation Graph}
\begin{algorithmic}[1]
\Require Database $\mathcal{D}$, max connections $M$
\For{each vector $v_i$ in $\mathcal{D}$}
    \State $neighbors \gets \text{findKNearest}(v_i, M)$
    \For{each neighbor $v_j$ in $neighbors$}
        \State $\text{graph.addEdge}(i, j)$
        \State $\text{graph.addEdge}(j, i)$ \Comment{Bidirectional}
    \EndFor
\EndFor
\end{algorithmic}
\end{algorithm}

\textbf{Complexity}: $O(n^2 d)$ for construction (can be optimized)

\subsubsection{Greedy Search Algorithm}
\begin{algorithm}
\caption{Graph-Based ANN Search}
\begin{algorithmic}[1]
\Require Query $q$, graph $G$, entry point $s$, $k$, search budget $L$
\Ensure Approximate top-$k$ results
\State Initialize visited set $V \gets \emptyset$
\State Initialize candidates heap $C$ (max heap by distance)
\State Initialize results heap $R$ (max heap by distance)
\State $C.\text{insert}((- \text{sim}(q, s), s))$
\State $R.\text{insert}((- \text{sim}(q, s), s))$
\While{$|V| < L$ \textbf{and} $C$ not empty}
    \State $(dist, u) \gets C.\text{extractMax}()$
    \If{$-dist < -R.\text{max}().dist$ \textbf{and} $|R| \geq k$}
        \State \textbf{break} \Comment{Early termination}
    \EndIf
    \State $V \gets V \cup \{u\}$
    \For{each neighbor $v$ of $u$ in $G$}
        \If{$v \notin V$}
            \State $sim \gets \text{similarity}(q, v)$
            \If{$sim > -R.\text{max}().dist$ \textbf{or} $|R| < k$}
                \State $C.\text{insert}((-sim, v))$
                \State $R.\text{insert}((-sim, v))$
                \If{$|R| > k$}
                    \State $R.\text{extractMax}()$
                \EndIf
            \EndIf
        \EndIf
    \EndFor
\EndWhile
\State \Return top $k$ from $R$
\end{algorithmic}
\end{algorithm}

\subsubsection{Complexity Analysis}
\begin{itemize}
    \item Time: $O(Ld)$ where $L$ is search budget
    \item Space: $O(k + L)$
    \item Typical speedup: $10-100\times$ over brute force
    \item Typical recall: $90-95\%$ with proper parameters
\end{itemize}

\subsubsection{Parameters}
\begin{itemize}
    \item \textbf{M}: Max connections per node (typically 5-16)
    \item \textbf{efConstruction}: Size of dynamic candidate list during construction
    \item \textbf{efSearch}: Size of dynamic candidate list during search (controls recall-speed tradeoff)
\end{itemize}

\subsection{KD-Tree Based Search}

\subsubsection{Structure}
KD-Tree (k-dimensional tree) is a space-partitioning data structure for organizing points in k-dimensional space. Each node represents a hyperplane that divides the space into two half-spaces.

\subsubsection{Construction Algorithm}
\begin{algorithm}
\caption{Build KD-Tree}
\begin{algorithmic}[1]
\Require Database $\mathcal{D}$, current dimension $depth$
\Ensure KD-Tree root node
\If{$\mathcal{D}$ is empty}
    \State \Return null
\EndIf
\State $axis \gets depth \bmod d$
\State Sort $\mathcal{D}$ by $axis$-th dimension
\State $median \gets \mathcal{D}[|\mathcal{D}|/2]$
\State Create node with $median$
\State $node.left \gets \text{buildKDTree}(\mathcal{D}[0:median], depth+1)$
\State $node.right \gets \text{buildKDTree}(\mathcal{D}[median+1:end], depth+1)$
\State \Return node
\end{algorithmic}
\end{algorithm}

\textbf{Complexity}: $O(dn \log n)$ for construction

\subsubsection{Search Algorithm}
\begin{algorithm}
\caption{KD-Tree k-NN Search}
\begin{algorithmic}[1]
\Require Query $q$, tree node $node$, $k$, best list $B$
\Ensure Top-$k$ nearest neighbors
\If{$node$ is null}
    \State \Return
\EndIf
\State $dist \gets \text{distance}(q, node.point)$
\If{$|B| < k$ \textbf{or} $dist < B.worst()$}
    \State Insert $(dist, node.point)$ into $B$
    \If{$|B| > k$}
        \State Remove worst from $B$
    \EndIf
\EndIf
\State $axis \gets node.depth \bmod d$
\If{$q[axis] < node.point[axis]$}
    \State Search $node.left$ first, then $node.right$ if needed
\Else
    \State Search $node.right$ first, then $node.left$ if needed
\EndIf
\end{algorithmic}
\end{algorithm}

\subsubsection{Complexity Analysis}
\begin{itemize}
    \item Time: $O(d \log n)$ average case, $O(dn)$ worst case
    \item Space: $O(n)$ for tree structure
    \item Performance degrades in high dimensions ($d > 20$)
    \item Effective for low to moderate dimensions
\end{itemize}

\subsubsection{Advantages and Limitations}
\begin{itemize}
    \item \textbf{Advantages}: Exact results, good for low dimensions, simple implementation
    \item \textbf{Limitations}: Curse of dimensionality, poor performance when $d > 20$
\end{itemize}

\subsection{Ball Tree Based Search}

\subsubsection{Structure}
Ball Tree recursively partitions data into nested hyperspheres (balls). Each node contains a centroid and radius that bounds all points in its subtree.

\subsubsection{Construction Algorithm}
\begin{algorithm}
\caption{Build Ball Tree}
\begin{algorithmic}[1]
\Require Database $\mathcal{D}$
\Ensure Ball Tree root node
\If{$|\mathcal{D}| \leq$ leaf\_size}
    \State \Return leaf node with $\mathcal{D}$
\EndIf
\State $centroid \gets \text{mean}(\mathcal{D})$
\State Find dimension with largest spread
\State Split $\mathcal{D}$ along this dimension at median
\State $radius \gets \max_{v \in \mathcal{D}} \text{dist}(v, centroid)$
\State Create node with $centroid$ and $radius$
\State $node.left \gets \text{buildBallTree}(\mathcal{D}_{left})$
\State $node.right \gets \text{buildBallTree}(\mathcal{D}_{right})$
\State \Return node
\end{algorithmic}
\end{algorithm}

\textbf{Complexity}: $O(dn \log n)$ for construction

\subsubsection{Search Algorithm}
\begin{algorithm}
\caption{Ball Tree k-NN Search}
\begin{algorithmic}[1]
\Require Query $q$, tree node $node$, $k$, best list $B$
\Ensure Top-$k$ nearest neighbors
\If{$node$ is leaf}
    \For{each point $p$ in $node$}
        \State Update $B$ with $(\text{dist}(q, p), p)$ if better
    \EndFor
    \State \Return
\EndIf
\State $dist\_to\_centroid \gets \text{distance}(q, node.centroid)$
\If{$dist\_to\_centroid - node.radius > B.worst()$ \textbf{and} $|B| = k$}
    \State \Return \Comment{Prune this subtree}
\EndIf
\State Recursively search closer child first
\State Recursively search farther child if needed
\end{algorithmic}
\end{algorithm}

\subsubsection{Complexity Analysis}
\begin{itemize}
    \item Time: $O(d \log n)$ average case, better than KD-Tree in high dimensions
    \item Space: $O(n)$ for tree structure
    \item More robust to high dimensions than KD-Tree
    \item Triangle inequality used for efficient pruning
\end{itemize}

\subsubsection{Advantages and Limitations}
\begin{itemize}
    \item \textbf{Advantages}: Better than KD-Tree for high dimensions, exact results, efficient pruning
    \item \textbf{Limitations}: Still affected by curse of dimensionality, construction overhead
\end{itemize}

\subsection{Embedding Generation}

\subsubsection{API Integration}
The system integrates with Google's Gemini API to generate embeddings:

\begin{lstlisting}
Vector<float> getEmbedding(string text) {
    Send HTTP POST request to Gemini API
    Parse JSON response
    Extract embedding values
    Return 128-dimensional vector
}
\end{lstlisting}

\subsubsection{Caching Strategy}
\begin{itemize}
    \item Cache frequently requested embeddings using HashMap
    \item Reduces API calls and improves response time
    \item Configurable cache size with LRU eviction
\end{itemize}

\section{Implementation Details}

\subsection{VectorDB Class}

\subsubsection{Interface}
\begin{lstlisting}
class VectorDB {
private:
    Vector<VectorItem*> vectors;
    int nextId;
public:
    VectorDB();
    ~VectorDB();
    int insert(const Vector<float>& embedding);
    VectorItem* get(int id);
    VectorItem* at(int index);
    int size() const;
};
\end{lstlisting}

\subsubsection{Memory Management}
\begin{itemize}
    \item Uses raw pointers to \texttt{VectorItem} for flexibility
    \item Destructor ensures proper cleanup to prevent leaks
    \item Copy constructor disabled to prevent accidental duplication
\end{itemize}

\subsection{SearchEngine Class}

\subsubsection{Interface}
\begin{lstlisting}
class SearchEngine {
private:
    VectorDB* db;
public:
    SearchEngine(VectorDB* database);
    Vector<Pair<int,float>> search(
        const Vector<float>& query, int k);
};
\end{lstlisting}

\subsubsection{Search Flow}
\begin{enumerate}
    \item Initialize MinHeap of size $k$
    \item Iterate through all vectors in database
    \item Compute cosine similarity for each vector
    \item Maintain top-k in heap
    \item Extract and reverse results
\end{enumerate}

\subsection{GraphANN Class}

\subsubsection{Interface}
\begin{lstlisting}
class GraphANN {
private:
    VectorDB* db;
    Graph<int> graph;
    int maxConnections;
    Vector<Pair<float, int>> findKNearest(int nodeId, int k);
public:
    GraphANN(VectorDB* database, int maxConn = 5);
    void buildGraph();
    Vector<Pair<int, float>> search(
        const Vector<float>& query, int k, int searchBudget);
};
\end{lstlisting}

\subsubsection{Graph Construction}
\begin{itemize}
    \item For each vector, find M nearest neighbors using brute force
    \item Create bidirectional edges
    \item Resulting graph has average degree $2M$
\end{itemize}

\subsubsection{Entry Point Selection}
\begin{itemize}
    \item Random selection (simple)
    \item Centroid-based (better quality)
    \item Multiple entry points (improved recall)
\end{itemize}

\subsection{Persistent Storage}

\subsubsection{File Format}
\begin{verbatim}
[Header]
    Magic Number: 0x56444221 (VDB!)
    Version: 1
    Vector Count: n
    Dimension: d
[B-Tree Index]
    Root offset
    Node size
    Tree structure
[Compressed Vectors]
    Vector 0: [Huffman encoded]
    Vector 1: [Huffman encoded]
    ...
    Vector n-1: [Huffman encoded]
\end{verbatim}

\subsubsection{Save Operation}
\begin{enumerate}
    \item Write header with metadata
    \item Build Huffman tree from all vectors
    \item Write Huffman tree structure
    \item Compress and write each vector
    \item Build B-Tree index mapping IDs to offsets
    \item Write B-Tree structure
\end{enumerate}

\subsubsection{Load Operation}
\begin{enumerate}
    \item Read and validate header
    \item Load Huffman tree
    \item Load B-Tree index
    \item Vectors loaded lazily on demand (or eagerly based on configuration)
\end{enumerate}

\section{Testing and Validation}

\subsection{Unit Tests}

\subsubsection{Data Structure Tests}
\begin{itemize}
    \item \texttt{test\_vector.cpp}: Dynamic array operations
    \item \texttt{test\_heap.cpp}: Min/max heap correctness
    \item \texttt{test\_graph.cpp}: Graph operations and traversal
    \item \texttt{test\_hashmap.cpp}: Hash table operations
    \item \texttt{test\_btree.cpp}: B-tree insertion, deletion, search
    \item \texttt{test\_huffman.cpp}: Compression and decompression
    \item \texttt{test\_list.cpp}: Linked list operations
    \item \texttt{test\_pair.cpp}: Pair comparison and assignment
\end{itemize}

\subsubsection{Core Module Tests}
\begin{itemize}
    \item \texttt{test\_embedding.cpp}: API integration
    \item \texttt{test\_vectorDB.cpp}: Database operations
    \item \texttt{test\_searchengine.cpp}: Brute force search correctness
    \item \texttt{test\_graphann.cpp}: Approximate search quality
\end{itemize}

\subsection{Test Cases}

\subsubsection{Correctness Tests}
\begin{enumerate}
    \item \textbf{Self-similarity}: Query vector should match itself with similarity $\approx 1.0$
    \item \textbf{Ranking order}: Results must be sorted by decreasing similarity
    \item \textbf{K-limit}: Must return exactly $k$ results (or fewer if database smaller)
    \item \textbf{Empty database}: Gracefully handle queries on empty database
    \item \textbf{Single vector}: Correctly handle database with one vector
\end{enumerate}

\subsubsection{Performance Tests}
\begin{enumerate}
    \item \textbf{Scalability}: Measure query time for $n \in \{10^2, 10^3, 10^4, 10^5\}$
    \item \textbf{Dimensionality}: Test with $d \in \{32, 64, 128, 256, 512\}$
    \item \textbf{K-parameter}: Vary $k \in \{1, 5, 10, 50, 100\}$
\end{enumerate}

\subsubsection{Recall Measurement}
For approximate search:
\begin{equation}
\text{Recall@k} = \frac{|\text{ANN\_results} \cap \text{exact\_results}|}{k}
\end{equation}








Target: Recall@10 $> 0.90$ with speedup $> 10\times$

\subsection{Test Results}

\subsubsection{Sample Output}
\begin{verbatim}
====================================
Search Engine Test Suite (128D)
====================================

=== Test 1: Basic Search ===
Inserted 3 vectors, searching for first vector
Top 3 results:
ID: 0, Similarity: 1.0
ID: 1, Similarity: 0.342
ID: 2, Similarity: 0.187
PASS: Query vector found itself with highest similarity

=== Test 2: Top-K Retrieval ===
Inserted 7 vectors (5 similar + 2 random)
Top 3 results:
Rank 1 - ID: 0, Similarity: 1.0
Rank 2 - ID: 1, Similarity: 0.923
Rank 3 - ID: 2, Similarity: 0.876
PASS: Results ranked in descending order

=== Test 3: Large Database Search ===
Inserting 1000 vectors...
Search completed in 234.5 ms
Top 10 results:
Rank 1 - ID: 0, Similarity: 0.987
...
Found 9/10 intentionally similar vectors in top 10
PASS: Found majority of similar vectors
\end{verbatim}

\section{Performance Analysis}

\subsection{Time Complexity Summary}

\begin{table}[h]
\centering
\begin{tabular}{@{}lllll@{}}
\toprule
\textbf{Operation} & \textbf{Brute Force} & \textbf{Graph ANN} & \textbf{KD-Tree} & \textbf{Ball Tree} \\ \midrule
Insert             & $O(1)$               & $O(Md)$            & $O(d\log n)$     & $O(d\log n)$      \\
Build Index        & N/A                  & $O(n^2d)$          & $O(dn\log n)$    & $O(dn\log n)$     \\
Search (avg)       & $O(nd)$              & $O(Ld)$            & $O(d\log n)$     & $O(d\log n)$      \\
Search (worst)     & $O(nd)$              & $O(nd)$            & $O(dn)$          & $O(dn)$           \\
Memory             & $O(nd)$              & $O(nd + nM)$       & $O(n)$           & $O(n)$            \\ 
\textbf{Best For}  & Small $n$            & Large $n$, high $d$ & Low $d$ ($<$20) & Medium $d$        \\ \bottomrule
\end{tabular}
\caption{Complexity comparison of search methods}
\end{table}

Where:
\begin{itemize}
    \item $n$ = number of vectors
    \item $d$ = vector dimension
    \item $k$ = number of results
    \item $M$ = max connections per node
    \item $L$ = search budget
\end{itemize}

\subsection{Space Complexity}

\subsubsection{Memory Usage}
\begin{itemize}
    \item Base vectors: $4nd$ bytes (32-bit floats)
    \item Graph edges: $8nM$ bytes (bidirectional, 4 bytes per ID)
    \item B-Tree index: $O(n)$ bytes
    \item Total: $O(nd + nM)$ bytes
\end{itemize}

\subsubsection{Disk Usage}
\begin{itemize}
    \item Uncompressed: $4nd$ bytes
    \item Huffman compressed: $\approx 2-2.5nd$ bytes (30-50\% reduction)
\end{itemize}

\subsection{Benchmark Results}

\subsubsection{Query Latency vs Database Size}
\begin{table}[h]
\centering
\begin{tabular}{@{}lrr@{}}
\toprule
\textbf{Database Size} & \textbf{Brute Force (ms)} & \textbf{Graph ANN (ms)} \\ \midrule
100                    & 2.3                       & 1.8                     \\
1,000                  & 23.5                      & 4.2                     \\
10,000                 & 234.7                     & 12.6                    \\
100,000                & 2,341.2                   & 45.3                    \\ \bottomrule
\end{tabular}
\caption{Query latency comparison (d=128, k=10)}
\end{table}

\subsubsection{Speedup Factor}
Graph ANN achieves $15-50\times$ speedup over brute force for databases larger than 10,000 vectors.

\subsubsection{Recall Quality}
\begin{table}[h]
\centering
\begin{tabular}{@{}lrr@{}}
\toprule
\textbf{Search Budget} & \textbf{Recall@10} & \textbf{Query Time (ms)} \\ \midrule
50                     & 0.76               & 8.2                      \\
100                    & 0.89               & 12.6                     \\
200                    & 0.94               & 23.4                     \\
500                    & 0.97               & 51.7                     \\ \bottomrule
\end{tabular}
\caption{Recall vs speed tradeoff (n=10,000, d=128)}
\end{table}

\subsection{Compression Analysis}

\begin{table}[h]
\centering
\begin{tabular}{@{}lrr@{}}
\toprule
\textbf{Vector Count} & \textbf{Uncompressed (MB)} & \textbf{Compressed (MB)} \\ \midrule
1,000                 & 0.5                        & 0.3                      \\
10,000                & 5.1                        & 3.2                      \\
100,000               & 51.2                       & 32.7                     \\
1,000,000             & 512.0                      & 327.0                    \\ \bottomrule
\end{tabular}
\caption{Storage requirements (d=128, 32-bit floats)}
\end{table}

Huffman compression achieves approximately 36\% reduction in storage size.

\section{Advanced Features}

\subsection{Metadata Filtering}

\subsubsection{Concept}
Allow filtering search results based on vector metadata (e.g., timestamp, category, tags).

\subsubsection{Implementation}
\begin{itemize}
    \item Extend \texttt{VectorItem} to include metadata fields
    \item Use segment tree for range queries (timestamp filtering)
    \item Use hash map for categorical filtering
    \item Apply filters during search before inserting into result heap
\end{itemize}

\subsection{Batch Operations}

\subsubsection{Bulk Insert}
\begin{itemize}
    \item Accept array of vectors for insertion
    \item Optimize graph construction by computing all neighbors in batch
    \item Reduces overhead compared to individual inserts
\end{itemize}

\subsubsection{Batch Query}
\begin{itemize}
    \item Process multiple queries simultaneously
    \item Share computation of vector norms and graph traversals
    \item Improves throughput for high query rates
\end{itemize}

\subsection{Update and Deletion}

\subsubsection{Vector Update}
\begin{itemize}
    \item Recompute affected edges in graph
    \item Update B-Tree mapping if vector moves on disk
    \item Invalidate cached similarity computations
\end{itemize}

\subsubsection{Vector Deletion}
\begin{itemize}
    \item Mark vector as deleted (soft delete)
    \item Skip during search operations
    \item Periodically compact database to reclaim space
    \item Update graph by removing node and reconnecting neighbors
\end{itemize}

\subsection{Distributed Architecture}

\subsubsection{Sharding Strategy}
\begin{itemize}
    \item Partition vectors across multiple nodes
    \item Hash-based or range-based sharding
    \item Each shard maintains its own graph
\end{itemize}

\subsubsection{Query Routing}
\begin{itemize}
    \item Query all shards in parallel
    \item Merge top-k results from each shard
    \item Return global top-k
\end{itemize}

\section{Learning Outcomes}

\subsection{Data Structures}

\subsubsection{Fundamental Understanding}
\begin{itemize}
    \item Implementing dynamic arrays with proper memory management
    \item Understanding amortized analysis for resizing operations
    \item Binary heap properties and heapify operations
    \item Graph representations and traversal algorithms
    \item Hash table design with collision resolution
    \item B-Tree structure for disk-based indexing
    \item Huffman encoding for optimal compression
\end{itemize}

\subsubsection{Advanced Concepts}
\begin{itemize}
    \item Template programming for generic data structures
    \item Memory allocation and deallocation strategies
    \item Copy constructors and assignment operators
    \item RAII (Resource Acquisition Is Initialization) principles
\end{itemize}

\subsection{Algorithms}

\subsubsection{Search Algorithms}
\begin{itemize}
    \item Linear search with heap-based top-k maintenance
    \item Graph traversal (BFS/DFS) for navigation
    \item Greedy search in graph structures
    \item Priority queue operations
\end{itemize}

\subsubsection{Optimization Techniques}
\begin{itemize}
    \item Early termination conditions
    \item Pruning strategies in search spaces
    \item Trade-offs between accuracy and speed
    \item Approximation algorithms
\end{itemize}

\subsection{System Design}

\subsubsection{Modularity}
\begin{itemize}
    \item Separation of concerns across modules
    \item Clean interfaces between components
    \item Dependency management
    \item Code organization and project structure
\end{itemize}

\subsubsection{Integration}
\begin{itemize}
    \item External API integration
    \item File I/O and serialization
    \item Error handling and validation
    \item Testing strategies
\end{itemize}

\subsection{Performance Engineering}

\subsubsection{Profiling}
\begin{itemize}
    \item Identifying performance bottlenecks
    \item Time complexity analysis
    \item Space complexity analysis
    \item Memory profiling
\end{itemize}

\subsubsection{Optimization}
\begin{itemize}
    \item Algorithm selection based on use case
    \item Data structure choice for specific operations
    \item Cache efficiency considerations
    \item Compression trade-offs
\end{itemize}

\section{Future Enhancements}

\subsection{Short-term Improvements}

\subsubsection{Multi-threading}
\begin{itemize}
    \item Parallelize similarity computations
    \item Concurrent graph construction
    \item Thread-safe operations for concurrent queries
    \item Lock-free data structures where applicable
\end{itemize}

\subsubsection{Better Indexing}
\begin{itemize}
    \item Implement full HNSW with hierarchical layers
    \item Product quantization for reduced memory footprint
    \item Inverted file index (IVF) for partitioning
\end{itemize}

\subsubsection{Enhanced API}
\begin{itemize}
    \item RESTful HTTP API for remote access
    \item Support multiple embedding models
    \item Batch API endpoints
    \item Authentication and authorization
\end{itemize}

\subsection{Long-term Vision}

\subsubsection{GPU Acceleration}
\begin{itemize}
    \item CUDA implementation for similarity computation
    \item GPU-based graph traversal
    \item Mixed precision computation (FP16/FP32)
\end{itemize}

\subsubsection{Advanced Indexing}
\begin{itemize}
    \item Locality-sensitive hashing (LSH)
    \item Tree-based methods (Annoy, Ball Tree)
    \item Learning-to-hash neural approaches
\end{itemize}

\subsubsection{Scalability}
\begin{itemize}
    \item Distributed multi-node architecture
    \item Dynamic sharding and rebalancing
    \item Replication for fault tolerance
    \item Streaming ingestion pipeline
\end{itemize}

\subsubsection{Monitoring and Observability}
\begin{itemize}
    \item Real-time performance metrics
    \item Query latency histograms
    \item Resource utilization tracking
    \item Alerting for anomalies
\end{itemize}

\section{Conclusion}

\subsection{Summary}
VectorDB successfully demonstrates the practical application of fundamental data structures and algorithms in building a modern, high-performance vector database system. The project implements:

\begin{itemize}
    \item 11 custom data structures from scratch
    \item Multiple search strategies (exact and approximate)
    \item Persistent storage with compression
    \item External API integration
    \item Comprehensive testing suite
\end{itemize}

The system achieves sub-second query latency for databases with up to $10^5$ vectors while maintaining high recall rates through graph-based approximate search. The modular architecture enables easy extension and maintenance.

\subsection{Key Achievements}

\begin{enumerate}
    \item \textbf{Educational}: Deep understanding of DSA concepts through implementation
    \item \textbf{Practical}: Working system applicable to real-world problems
    \item \textbf{Performance}: Competitive performance with modern vector databases
    \item \textbf{Code Quality}: Clean, well-structured, tested codebase
\end{enumerate}

\subsection{Lessons Learned}

\subsubsection{Technical}
\begin{itemize}
    \item Memory management is critical in C++ systems
    \item Template programming enables code reuse
    \item Choosing the right data structure dramatically affects performance
    \item Approximation algorithms are essential for scalability
\end{itemize}

\subsubsection{Engineering}
\begin{itemize}
    \item Testing is crucial for correctness
    \item Modularity simplifies debugging and extension
    \item Documentation improves maintainability
    \item Performance measurement guides optimization
\end{itemize}

\subsection{Impact}
This project bridges the gap between theoretical DSA knowledge and practical system building. It demonstrates that understanding fundamentals enables creation of sophisticated systems that power modern AI applications.

\subsection{Final Thoughts}
Vector databases represent a critical component of the AI infrastructure stack. As language models and embedding techniques continue to advance, efficient vector search will become increasingly important. This project provides a solid foundation for understanding both the challenges and solutions in this domain, preparing the way for more advanced research and development in high-dimensional data management.

\section{References}

\begin{enumerate}
    \item Malkov, Y. A., \& Yashunin, D. A. (2018). \textit{Efficient and robust approximate nearest neighbor search using Hierarchical Navigable Small World graphs}. IEEE transactions on pattern analysis and machine intelligence, 42(4), 824-836.
    
    \item Johnson, J., Douze, M., \& Jégou, H. (2019). \textit{Billion-scale similarity search with gpus}. IEEE Transactions on Big Data, 7(3), 535-547.
    
    \item Andoni, A., \& Indyk, P. (2008). \textit{Near-optimal hashing algorithms for approximate nearest neighbor in high dimensions}. Communications of the ACM, 51(1), 117-122.
    
    \item Cormen, T. H., Leiserson, C. E., Rivest, R. L., \& Stein, C. (2009). \textit{Introduction to algorithms} (3rd ed.). MIT press.
    
    \item Goodfellow, I., Bengio, Y., \& Courville, A. (2016). \textit{Deep learning}. MIT press.
    
    \item Mikolov, T., Chen, K., Corrado, G., \& Dean, J. (2013). \textit{Efficient estimation of word representations in vector space}. arXiv preprint arXiv:1301.3781.
    
    \item Jégou, H., Douze, M., \& Schmid, C. (2010). \textit{Product quantization for nearest neighbor search}. IEEE transactions on pattern analysis and machine intelligence, 33(1), 117-128.
    
    \item Huffman, D. A. (1952). \textit{A method for the construction of minimum-redundancy codes}. Proceedings of the IRE, 40(9), 1098-1101.
    
    \item Bayer, R., \& McCreight, E. (1972). \textit{Organization and maintenance of large ordered indices}. Acta informatica, 1(3), 173-189.
    
    \item Google AI. (2024). \textit{Gemini API Documentation}. Retrieved from https://ai.google.dev/
    
    \item Pinecone Systems. (2024). \textit{Vector Database Technology Overview}. Retrieved from https://www.pinecone.io/
    
    \item Weaviate. (2024). \textit{Vector Search Engine Documentation}. Retrieved from https://weaviate.io/
    
    \item Aumüller, M., Bernhardsson, E., \& Faithfull, A. (2020). \textit{ANN-Benchmarks: A benchmarking tool for approximate nearest neighbor algorithms}. Information Systems, 87, 101374.
    
    \item Li, W., Zhang, Y., Sun, Y., Wang, W., Li, M., Zhang, W., \& Lin, X. (2020). \textit{Approximate nearest neighbor search on high dimensional data—experiments, analyses, and improvement}. IEEE Transactions on Knowledge and Data Engineering, 32(8), 1475-1488.
    
    \item Stroustrup, B. (2013). \textit{The C++ programming language} (4th ed.). Addison-Wesley.
\end{enumerate}

\newpage

\section*{Appendix A: Code Listings}

\subsection*{A.1 Vector Class Implementation}
\begin{lstlisting}
template <typename T>
class Vector {
    T* arr;
    int size;
    int capacity;
public:
    Vector() {
        arr = new T[1];
        size = 0;
        capacity = 1;
    }
    
    void push_back(const T& val) {
        if(size == capacity) {
            capacity *= 2;
            T* temp = new T[capacity];
            for(int i = 0; i < size; i++)
                temp[i] = arr[i];
            delete[] arr;
            arr = temp;
        }
        arr[size++] = val;
    }
    
    T& operator[](int index) {
        return arr[index];
    }
    
    int length() const {
        return size;
    }
};
\end{lstlisting}

\subsection*{A.2 Similarity Metrics}
\begin{lstlisting}
class Similarity {
public:
    static float cosineSimilarity(
        const Vector<float>& a, 
        const Vector<float>& b) {
        float dot = 0, magA = 0, magB = 0;
        for(int i = 0; i < a.length(); i++) {
            dot += a[i] * b[i];
            magA += a[i] * a[i];
            magB += b[i] * b[i];
        }
        return dot / (sqrt(magA) * sqrt(magB));
    }
};
\end{lstlisting}

\section*{Appendix B: Performance Graphs}

\subsection*{B.1 Query Latency vs Database Size}
[Placeholder for graph showing logarithmic performance improvement with Graph ANN]

\subsection*{B.2 Recall vs Search Budget}
[Placeholder for graph showing recall improving with larger search budgets]

\subsection*{B.3 Memory Usage}
[Placeholder for graph comparing uncompressed vs compressed storage]

\section*{Appendix C: Project Structure}

\begin{tcolorbox}[
    title=Project Architecture (as implemented),
    colback=gray!5,
    colframe=blue!75!black,
    fonttitle=\bfseries,
    arc=0mm
]
\dirtree{%
.1 VectorDB/.
.2 backend/.
.3 core/.
.4 Embedding/.
.5 Embedding.h.
.4 vector/.
.5 VectorItem.h.
.5 VectorDB.h.
.5 Similarity.h.
.5 SearchEngine.h.
.5 GraphANN.h.
.5 KDANN.h.
.5 BallANN.h.
.4 storage/.
.5 HuffmanCompression.h.
.5 BtreeIndex.h.
.3 dataStructures/.
.4 Vector.h.
.4 Heap.h.
.4 Graph.h.
.4 HashMap.h.
.4 BTree.h.
.4 Huffman.h.
.4 List.h.
.4 Queue.h.
.4 Stack.h.
.4 Pair.h.
.4 AVL.h.
.3 common/.
.4 EnvUtil.h.
.3 external/.
.4 nlohmann/json.hpp.
.3 EmbeddingCache.h.
.3 main.cpp.
.3 tests/.
.4 test\_vector.cpp, test\_heap.cpp, test\_hashmap.cpp, \ldots\ (one per data structure and core class).
.4 benchmark\_cpp.cpp.
.4 benchmark\_recall.cpp.
.3 docs/.
.4 project\_documentation.tex.
.4 project\_documentation.pdf.
.1 CMakeLists.txt.
.1 run\_tests.sh.
.1 .github/workflows/ci.yml.
.1 README.md.
.1 build-plan.md.
}
\end{tcolorbox}

\section*{Appendix D: Implementation Status}

This appendix was added after the initial submission to keep the report honest about what is actually built versus what was originally aspirational. The rest of this document was written describing the intended design; as of this appendix, the implementation has largely caught up to it:

\begin{itemize}
    \item \textbf{Implemented and tested}: all 11 data structures (including AVL, not in the original count), \texttt{VectorDB}, \texttt{SearchEngine}, \texttt{GraphANN}, \texttt{KDANN}, \texttt{BallANN}, the Gemini embedding client, an in-memory LRU embedding cache, soft delete, vector update, single-key metadata filtering, B-tree + Huffman backed persistence (\texttt{VectorDB::save}/\texttt{load}), and a minimal CLI (\texttt{backend/main.cpp}) supporting seed-and-query, \texttt{save}, \texttt{load}, and \texttt{insert}.
    \item \textbf{Persistence, honestly scoped}: vectors are Huffman-compressed as one shared block (one tree, better compression, matches how \texttt{Huffman.h}'s \texttt{encode}/\texttt{decode} are actually meant to be used) rather than per-vector blocks with independent byte offsets. The on-disk ``B-tree index'' is therefore a flat id table, not a serialized tree with pointers; the in-memory \texttt{BtreeIndex} built from it \emph{is} a real \texttt{BTree}, and is what every \texttt{VectorDB::get()} call goes through. Loading is eager (whole file decompressed at once) rather than lazy. Per-item metadata is not currently persisted.
    \item \textbf{Not implemented}: an HTTP API layer, multi-threading, product quantization / IVF indexing, and the distributed/sharding architecture described in the Advanced Features and Future Enhancements sections above — all still aspirational.
    \item \textbf{On the Performance Analysis numbers earlier in this report}: the query latency, recall, and compression tables in Section 8 were originally illustrative targets, not measurements. Real measurements now exist: \texttt{backend/tests/benchmark\_cpp.cpp} runs the full pipeline against 95 live-embedded words; \texttt{backend/tests/benchmark\_recall.cpp} measures Recall@10 and speedup for \texttt{GraphANN} against \texttt{SearchEngine} on synthetic data from $n=1{,}000$ to $8{,}000$. At that scale, with \texttt{maxConnections=10} and a fixed search budget of $2k$, measured speedup reached roughly $4\times$ (not the $15$--$50\times$ originally claimed) and recall \emph{declined} from about $0.80$ to $0.55$ as $n$ grew — because the search budget is a constant, not scaled to database size. This is a legitimate, useful finding, not a bug: the original targets assumed a tunable \texttt{efSearch} parameter (Section 5.3.4) that the current \texttt{GraphANN} does not yet expose. Widening the search budget as $n$ grows is the natural next step and is tracked in \texttt{build-plan.md}.
\end{itemize}

\end{document}
